\section{Integral primitive classes and stable representatives}
\label{sec:generators}

Fix an odd prime $p$.  Write $G_a=\GL_a(\Zp)$, the group denoted
$K_a$ in Section~\ref{sec:constants}, and let $S$ be the
$p$-complete sphere.  The continuous cochain spectrum
$C^*_{\mathrm{cts}}(G_a;E)$ is the one defined in
Section~\ref{sec:constants}: for bounded finite $p$-primary Postnikov
coefficients it is continuous totalization, and complete coefficients
are obtained by the compatible Postnikov and $p$-completion limits.
In particular, its cohomological filtration has spectral sequence
\begin{equation}\label{gen:ahss}
 E_2^{s,t}=H^s_{\mathrm{cts}}(G_a,\pi_tE)
 \quad\Longrightarrow\quad
 \pi_{t-s}C^*_{\mathrm{cts}}(G_a;E).
\end{equation}
For connective $E$, the coefficient map $E\to H\pi_0E$ sends a
class of degree $-b$ to a class in
$H^b_{\mathrm{cts}}(G_a,\pi_0E)$.  We call this its leading ordinary
class; it is defined even when it is zero.  The source class has
filtration at least $b$, and this coefficient image is its component
in filtration $b$.  For $\pi_0E=\Zp$, reduction modulo $p$ gives its
ordinary mod-$p$ Hurewicz class.

Theorem~\ref{thm:constant-cohomology} supplies the following facts for
$a\leq p-2$.  The group $G_a$ has $p$-cohomological dimension $a^2$;
its integral constant cohomology is finitely generated and
$p$-torsion-free; and its mod-$p$ cohomology has the graded dimensions
of the exterior algebra on degrees $1,3,\ldots,2a-1$.
We now construct compatible integral classes and lift them to sphere
cochains.  The regulator and top-volume comparisons then show that
these selected classes generate the exterior algebras integrally and
modulo $p$.  Their normalization and compatibility as the rank varies
will be needed in the comparison with the full frame tower.

\begin{theorem}[Stable primitive generators]\label{thm:stable-generators}
Let $1\leq N\leq p-2$.  There are classes
\[
 c_{a,i}\in
 \pi_{1-2i}C^*_{\mathrm{cts}}(G_a;S),
 \qquad 1\leq i\leq a\leq N,
\]
with leading integral classes
\[
 \alpha_{a,i}\in H^{2i-1}_{\mathrm{cts}}(G_a,\Zp),
\]
having the following properties.
\begin{enumerate}
\item The classes $\alpha_{a,i}$ induce isomorphisms of graded algebras
\[
 \bigwedge_{\Zp}(\alpha_{a,1},\ldots,\alpha_{a,a})
 \xrightarrow{\ \sim\ } H^*_{\mathrm{cts}}(G_a,\Zp),
 \qquad
 \bigwedge_{\Fp}(\overline\alpha_{a,1},\ldots,
                  \overline\alpha_{a,a})
 \xrightarrow{\ \sim\ } H^*_{\mathrm{cts}}(G_a,\Fp).
\]
Their rationalizations are nonzero multiples of the standard primitive
classes.  For $i=1$ the class is the divided determinant logarithm.
\item For the upper-left block inclusion $G_b\longrightarrow G_a$,
with $b\leq a$, one has
\[
 \operatorname{res}\alpha_{a,i}=\alpha_{b,i}\quad(i\leq b),
 \qquad
 \operatorname{res}\alpha_{a,i}=0\quad(b<i\leq a).
\]
The sphere representatives can be chosen so that
$\operatorname{res}c_{a,i}=c_{b,i}$ for $i\leq b$.
\item Apply Proposition~\ref{prop:modification-transport} at rank $N$.
The classes $c_{N,i}$ have transported representatives
\[
 u_i\in\pi_{1-2i}C^*_{\mathrm{cts}}(P_N;S).
\]
Their leading mod-$p$ classes, denoted $\xi_i$, are the transported
classes $\overline\alpha_{N,i}$.  On the extension space with a marked
connected quotient and an $m$-dimensional Tate kernel, their constant
cohomology action is the upper-left block restriction of
$\overline\alpha_{N,i}$, hence is given by
$\overline\alpha_{m,i}$ for $i\leq m$ and vanishes for $i>m$.
\end{enumerate}
\end{theorem}

The construction of the classes and their integrality occupy the rest
of this section.  The last assertion of the theorem is a compatibility
of the actual cochain maps with the modification comparison.  It does
not use the coefficient-module calculation of
Proposition~\ref{prop:constant-module}.

\subsection{The continuous class from algebraic K-theory}
\label{gen:continuous-k-section}

At each positive integer $r$, the standard representation of
$\GL_a(\Z/p^r)$ gives a based additive class with coefficients in
$K(\Z/p^r)$.  We use its reduced version, obtained by subtracting its
rank, and denote it by $\kappa_{a,r}$.  These classes are natural for
reduction in $r$.  Under block sum, their pullbacks add.  These
statements hold already for the maps to algebraic $K$-theory, before
taking homotopy groups.

\begin{lemma}[Continuous evaluation of the standard representation]
\label{gen:k-continuity}
The classes $\kappa_{a,r}$ define a continuous additive class
\[
 \kappa_a\in
 \pi_0 C^*_{\mathrm{cts}}(G_a;K(\Zp)^\wedge_p).
\]
If $f:K(\Zp)^\wedge_p\longrightarrow E$ is a map to a complete
connective spectrum of finite type, composition with $f$ is represented
in the continuous Postnikov cochain construction.  It is compatible
with restriction to blocks and with block sums.  Restriction to the
discrete group $G_a$ agrees with evaluation of the usual discrete
algebraic $K$-theory class followed by completion and $f$.
\end{lemma}

\begin{proof}
First fix a coefficient exponent $p^l$ and a bounded Postnikov
interval.  The $p$-adic continuity theorem applies to $\Zp$, since it
is local, $p$-torsion-free, and henselian along $(p)$.  It identifies
\[
 K(\Zp)/p^l\longrightarrow
 \{K(\Z/p^r)/p^l\}_r
\]
on the pro systems of homotopy groups in every fixed degree; see
\cite[Theorem~C]{GeisserHesselholt}.  The mod-$p$ formulation also appears in
\cite[Section~5.2, Theorem~5.10]{CMM}; the assertion for $p^l$ follows
from it by the finite coefficient exact sequences.
The relevant homotopy groups of the finite rings are finite.  In
positive degrees this is the usual finiteness of the $K$-groups of a
finite ring \cite[Chapter~IV, Proposition~1.16]{WeibelKBook}; it can also be obtained from stability for its finite
linear groups and rational acyclicity.  Degree zero becomes finite
after reduction modulo $p^l$.

Consequently the displayed map is a pro-Postnikov equivalence on the
chosen interval.  Explicitly, the inverse systems of finite groups
are Mittag--Leffler.  An isomorphism from a finite constant system to
their limit has pro-zero kernel and cokernel: after passing to stable
images, the system is an inverse system of surjections, and the finite
limit detects its stable image at every fixed stage.  Induction through
the finitely many Postnikov fibers gives the corresponding statement
for spectra on the interval.  Alternatively, the pro homotopy-group
form of the cited continuity theorem gives this directly.

For these bounded targets the natural maps out of the constant system
can therefore be computed after passage far enough down the finite-ring
tower.  The representative $\kappa_{a,r}$ then factors through the
finite congruence quotient of $G_a$ and is a continuous cochain
representative.  The pro equivalence retains the homotopies between
these representatives, so the construction is compatible as both the
Postnikov interval and $l$ increase.  Taking these complete limits
defines $\kappa_a$ and its composites with $f$.  The finite-stage
construction commutes with block sum, and therefore so do the limits.
Finally, every comparison was induced by reduction of the usual
standard representation, which proves the assertion about discrete
evaluation.

Only bounded Postnikov intervals are needed to form any fixed ordinary
cohomology component.  Thus this construction of the ordinary
components is available at arbitrary ranks, even when we will use a
finite cohomological dimension bound only in the range $a\leq p-2$.
\end{proof}

Let $\ell_p$ denote the connective Adams summand of $p$-complete complex
$K$-theory, so that
\[
 \pi_*\ell_p=\Zp[v],\qquad |v|=2p-2.
\]
We fix the choices in the following local calculation once for all
ranks.

\begin{lemma}[Local K-theory coordinates]
\label{gen:local-k-coordinates}
For $2\leq i\leq p-2$, there are spectrum maps
\[
 f_i:K(\Zp)^\wedge_p\longrightarrow\Sigma^{2i-1}\ell_p
\]
which induce a generator coordinate on the free group in degree
$2i-1$.  Via the integral etale edge and localization for $\Qp$, this
coordinate identifies that group with
$H^1(\Qp,\Zp(i))$, up to a $p$-adic unit.
\end{lemma}

\begin{proof}
Use the connective-cover comparison between local algebraic
$K$-theory and topological cyclic homology
\cite[Theorem~D and Addendum~5.2]{HesselholtMadsenWitt}, together with
the odd-prime decomposition of $\operatorname{TC}(\Z)^\wedge_p$; see
\cite[Sections~2 and~7]{BlumbergMandellTC}.  The summands indexed by twists
$i=2,\ldots,p-2$ are the connective suspensions
$\Sigma^{2i-1}\ell_p$.  Projection to one of them gives $f_i$.
Only these nonexceptional twist residues are used.  In particular the
argument uses the local calculation and imposes no regular-prime
condition on global algebraic $K$-theory.

Put $d=2i-1$ and $X=K(\Zp)^\wedge_p$.  The same decomposition gives
\[
 \pi_dX=\Zp,\qquad \pi_{d-1}X=0,\qquad
 g_i:=(f_i)_*:\pi_dX\xrightarrow{\sim}\pi_0\ell_p=\Zp.
\]
To check the adjacent degree, write $L$ for the periodic $p$-complete
Adams summand and $J=L_{K(1)}S$.  The decomposition has summands
$j$, $\Sigma j'$, and $z_0,\ldots,z_{p-2}$.  Here $j$ and $j'$ are
the connective covers of $J$ and $J'$, with $J'$ noncanonically
equivalent to $J$; $z_k$ is the $1$-connected cover of
$\Sigma^{2k-1}L$ for $k\ne0$, and $z_0$ is the $(-2)$-connected
cover of $\Sigma^{-1}L$.  Passing to $X$ takes the connective cover
of this wedge.  Now $\pi_*L$ is concentrated in degrees divisible
by $2(p-1)$, while $\pi_*J$ is concentrated in degree zero and in
degrees congruent to $-1$ modulo $2(p-1)$.  Thus, in the range
$3\leq d\leq2p-5$, only the bottom group of the summand
$\Sigma^d\ell_p$ contributes in degree $d$, and no summand contributes
in degree $d-1$.  These descriptions are in
\cite[Section~2]{BlumbergMandellTC}.

Write $K_j(A;\Z/p^r)=\pi_j(K(A)/p^r)$.  For a connective spectrum,
the completion map is an equivalence modulo $p^r$: it is an
equivalence modulo $p$, and the Moore spectrum for $p^r$ is built by
finitely many extensions from the Moore spectrum for $p$.  The
coefficient exact sequence and $\pi_{d-1}X=0$ therefore identify
\[
 (\pi_dX)/p^r\xrightarrow{\sim}K_d(\Zp;\Z/p^r).
\]
The localization sequence for $\Zp\subset\Qp$ has residue term
$K(\Fp)$.  By \cite[Theorem~B and the following paragraph]
{HesselholtMadsenWitt}, $K(\Fp)^\wedge_p\simeq H\Zp$, so
$K(\Fp)/p^r\simeq H(\Z/p^r)$.  Since $d\geq3$, localization and
\cite[Theorem~A, the descent spectral sequence~(6.1.9), and
Theorem~6.1.10]{HesselholtMadsenLocalFields} consequently give
compatible isomorphisms
\[
 (\pi_dX)/p^r
 \xrightarrow{\sim}K_d(\Qp;\Z/p^r)
 \xrightarrow{\sim}H^1(\Qp,\mu_{p^r}^{\otimes i}).
\]
The last theorem applies to the complete characteristic-zero
discrete valuation field $\Qp$, whose residue field is perfect of
characteristic $p>2$.  Its odd-degree isomorphism is the unique
$H^1$ edge of the canonical etale comparison and is natural in the
coefficient maps.  Continuous integral cochains are the inverse limit
of the finite-coefficient cochains.  Their transitions are surjective
by choosing set sections of the finite coefficient maps, and the
finite groups in degree zero have vanishing $\varprojlim^1$.
Taking the inverse limit gives an isomorphism
\[
 \operatorname{edge}_i:\pi_dX\xrightarrow{\sim}
 H^1(\Qp,\Zp(i)).
\]
Consequently $g_i\circ\operatorname{edge}_i^{-1}$ is an integral
isomorphism from this etale group to $\Zp$.  The universal Chern
regulator used below is compared with these coordinates in
Lemma~\ref{gen:soule-coordinate}.
\end{proof}

Compose $\kappa_a$ with $f_i$, and write the resulting class as
\begin{equation}\label{gen:ell-class}
 z_{a,i}\in
 \pi_{1-2i}C^*_{\mathrm{cts}}(G_a;\ell_p).
\end{equation}
Its leading integral class is denoted $\alpha_{a,i}$.  These leading
classes are defined at every rank by the bounded construction in
Lemma~\ref{gen:k-continuity}.

\subsection{Rational primitivity and the lift through the image of J}
\label{gen:sphere-lifts-section}

Write $b_i$ for the standard rational primitive in degree $2i-1$ in
matrix Lie algebra cohomology.  We normalize it as the transgression
of the algebraic de Rham Chern class $c_i$.  Equivalently it is the
primitive used to define the $p$-adic Borel regulator in
\cite[Definition~0.4.5 and Remark~0.4.6]{HuberKings}.  Restriction to
$\mathfrak{gl}_a$ is zero for $i>a$.  The same notation will be used
for its image in continuous rational group cohomology under the
Lazard comparison.

\begin{lemma}[Rational components of the additive class]
\label{gen:rational-components}
For $a\leq p-2$ and $i\leq a$, the only possibly nonzero rational
ordinary-cohomology component of $z_{a,i}$ is its component in degree
$2i-1$.  Moreover $\alpha_{a,i}\otimes\Qp$ is a scalar multiple of
$b_i$.
\end{lemma}

\begin{proof}
The standard representation is additive for block sum and each $f_i$
is a map of spectra.  Thus every rational ordinary component of
$z_{a,i}$ is primitive under block sum.  To interpret this statement
stably, work first in a fixed bounded Postnikov interval of the
coefficient spectrum.  Lemma~\ref{gen:k-continuity} constructs its
components at all matrix ranks, compatibly.  Clear the finitely many
denominators on this interval and then rationalize the resulting
continuous integral cohomology classes.  This gives the usual stable
primitive condition without any exchange of an unbounded Postnikov
limit with rationalization.

The primitive space in stable rational matrix Lie cohomology is
one-dimensional in each degree $2j-1$, generated by $b_j$, and its
restriction to rank $a$ vanishes for $j>a$.  The possible ordinary
components of \eqref{gen:ell-class} have degrees
\[
 2i-1+2r(p-1)=2\bigl(i+r(p-1)\bigr)-1,
 \qquad r\geq0.
\]
If $r>0$, the corresponding primitive index exceeds $a$, since
$a\leq p-2$.  These components therefore vanish.  The component for
$r=0$ is a multiple of $b_i$, as asserted.
\end{proof}

\begin{lemma}[Sphere representatives]\label{gen:sphere-representatives}
For $1\leq i\leq a\leq p-2$, the class $\alpha_{a,i}$ has a
representative in
$\pi_{1-2i}C^*_{\mathrm{cts}}(G_a;S)$.
For $i\geq2$ its image in $\ell_p$ is $z_{a,i}$.
For $i=1$ one can take the pullback of the divided logarithm along
the determinant.
\end{lemma}

\begin{proof}
Put $q_0=2p-2$, and choose $q\in\Zp^\times$ whose Adams operation
generates the pro-Adams action and whose finite residue is primitive.
The operation $\psi^q-1$ on $\ell_p$ is zero on $\pi_0$, so it factors
through the positive connective cover, identified with
$\Sigma^{q_0}\ell_p$.  Denote the resulting map by
\[
 \theta:\ell_p\longrightarrow\Sigma^{q_0}\ell_p,
 \qquad j_p=\fib(\theta).
\]
This is the connective image-of-$J$ fiber, and the sphere unit lifts
to a map $S\longrightarrow j_p$.

For any fixed suspension of $\ell_p$, spectral sequence
\eqref{gen:ahss} has finitely many contributing groups on each total
degree, since $s\leq a^2$.  All of them are finite free over $\Zp$ by
Theorem~\ref{thm:constant-cohomology}.  The rational spectral sequence
has zero differentials, because the rational coefficient spectrum is
a direct sum of ordinary Eilenberg--Mac Lane spectra in this finite
range.  Induction on the pages shows that the integral differentials
also vanish: a map between the free groups on a page which vanishes
rationally is zero.  The resulting finite extensions of free
$\Zp$-modules are free.  In particular the abutment in each fixed
total degree is torsion-free and injects into its rationalization.

By Lemma~\ref{gen:rational-components}, the rational class $z_{a,i}$
has only its degree-zero coefficient component.  The map $\theta$
kills this component.  Thus $\theta z_{a,i}$ is rationally zero, and
the torsion-freeness just proved makes it zero integrally.  A choice
of its nullhomotopy lifts $z_{a,i}$ to cochains with coefficients in
$j_p$.

The classical odd-primary image-of-$J$ range identifies $S\to j_p$
on homotopy below the first beta stem
\[
 b_p=2p(p-1)-2.
\]
We use only degrees at most $a^2$, and
\begin{equation}\label{gen:beta-range}
 a^2\leq(p-2)^2<2p(p-1)-2.
\end{equation}
The finite continuous cohomological dimension bound shows that
coefficients above degree $a^2$ cannot affect a negative-degree
cochain class.  Indeed a coefficient in degree $t>a^2$ contributes
only total degrees $t-s>0$.  Therefore $S\to j_p$ induces the needed
isomorphism in degree $1-2i$ on continuous cochains, and the lifted
$j_p$ class has a sphere representative.  We use here the usual first
image-of-$J$ range, as recorded in the classical calculation of the
odd-primary stable stems; see
\cite[Theorems~1.1.14 and~1.5.19(a)]{RavenelGreenBook}.

For $i=1$, normalize the logarithm by
\[
 \lambda:\Zp^\times\longrightarrow\Zp,
 \qquad \lambda(u)=\frac{\log u}{\log(1+p)}.
\]
It kills the Teichmuller factor and sends $1+p$ to $1$.  The group
$\Zp^\times$ has $p$-cohomological dimension one.  In total degree
$-1$, its connective sphere-cochain spectral sequence therefore has
just the entry $H^1(\Zp^\times,\Zp)$.  Thus $\lambda$ lifts to a
sphere cochain, whose pullback along $\det:G_a\to\Zp^\times$
defines $c_{a,1}$ and $\alpha_{a,1}$.
\end{proof}

\subsection{The integral regulator lattice}
\label{gen:regulator-section}

The preceding argument does not yet show that the leading classes are
generators modulo $p$.  For that purpose we first determine their
rational normalization relative to $b_i$.

We first record the compact-coefficient finiteness used below.
Let $T=\Zp(i)$, $T_r=T/p^rT$, and
$M=H^1_{\mathrm{cts}}(G_{\Qp},T)$.
Finite-coefficient local Galois cohomology is finite
\cite[Chapter~I, Theorem~2.1]{MilneADT}.
Continuous cochains with coefficients in $T$ are the inverse limit
of the $T_r$ cochains, whose transitions are surjective by choosing
set sections of the finite coefficient maps.  The finite groups
$H^0(G_{\Qp},T_r)$ have zero $\varprojlim^1$, so
$M=\varprojlim_r H^1(G_{\Qp},T_r)$ is a compact pro-$p$ module.
The coefficient sequence for multiplication by $p$ gives an injection
$M/pM\hookrightarrow H^1(G_{\Qp},T_1)$.
Choose lifts $m_1,\ldots,m_e$ of a finite basis of $M/pM$.
Successive reduction modulo $p$ expresses any element as a
$\Zp$-linear combination of these lifts plus a remainder in $p^NM$.
That remainder tends to zero: its projection to the group with
coefficients $T_r$ is zero once $N\geq r$.
The coefficient sums converge in $\Zp$, proving that $M$ is
finitely generated.  This supplies the compact integral assertion
without applying a theorem about discrete finitely generated
modules to $T$.

\begin{lemma}[Integral Soul\'e coordinate]\label{gen:soule-coordinate}
Let $2\leq i\leq p-2$, $d=2i-1$, and $X=K(\Zp)^\wedge_p$.
There is a $\Zp$-linear isomorphism
\[
 c_i^{\mathrm S}:\pi_dX\xrightarrow{\sim}H^1(\Qp,\Zp(i))
\]
whose value on the completion of $x\in K_d(\Zp)$ rationalizes to
the Soul\'e regulator $r_p(x)$ of \cite[Section~1.3]{HuberKings},
restricted from $\Qp$ by localization.
\end{lemma}

\begin{proof}
Write $M_i=H^1(\Qp,\Zp(i))$.  First match the coefficient conventions.
For a connective $K$-theory spectrum $Y$, the mod-$q$ Moore space
$P^d(q)$ used in \cite[Section~1, Example~1.4]{WeibelEtaleChern}
has suspension spectrum $\Sigma^{d-1}S/q$.  Adjunction and the
fiber/cofiber triangle give
\[
 [P^d(q),\Omega^\infty_0Y]
 =\pi_{d-1}F(S/q,Y)
 =\pi_{d-1}\Sigma^{-1}(Y/q)
 =\pi_d(Y/q).
\]
Here $d\geq3$, so the connected Moore space maps to the identity
component.  With the common cofiber orientation these identifications
commute with ordinary reduction and coefficient maps.

Compose the finite-coefficient identifications in
Lemma~\ref{gen:local-k-coordinates} with the universal etale Chern maps
\[
 c_{i,r}^{\mathrm S}:(\pi_dX)/p^r
 \xrightarrow{\sim}K_d(\Qp;\Z/p^r)
 \xrightarrow{c_{i,1}}H^1(\Qp,\mu_{p^r}^{\otimes i}).
\]
The last map is defined in \cite[Section~2, equation~(2.1)]
{WeibelEtaleChern} by the standard representation's universal Chern
class, its Kunneth component, and the mod-$p^r$ Hurewicz map.
Its coefficient prime is invertible in $\Qp$.  Proposition~2.4
there gives additivity for odd coefficients, and Proposition~2.8
gives compatibility with coefficient reductions for $d\geq2$.
Their inverse limit is therefore a continuous additive, hence
$\Zp$-linear, map $c_i^{\mathrm S}:\pi_dX\to M_i$.
The compatible etale edges identify the projection to coefficients
$\Z/p$ with the quotient $M_i\to M_i/pM_i$, and thus
$M_i/pM_i\cong H^1(\Qp,\mu_p^{\otimes i})$.

We prove that the mod-$p$ Chern map over $\Qp$ is onto.  Put
$L=\Qp(\zeta_p)$.  The polynomial $\Phi_p(1+T)$ is Eisenstein, so
$[L:\Qp]=p-1$.  By \cite[Proposition~5.2]{WeibelEtaleChern}, the
cokernel of
\[
 c_{i,1}:K_{2i-1}(L;\Z/p)\longrightarrow
 H^1(L,\mu_p^{\otimes i})
\]
is annihilated by $(i-1)!$.  Its hypotheses are $i\geq2$ and
$q\not\equiv2\pmod4$ for a field containing $1/q$ and a primitive
$q$th root.  They hold for $q=p$ and $L$ above.  Since
$(i-1)!$ is a unit modulo $p$, this map is onto.

For the finite separable extension $\Qp\subset L$,
\cite[Proposition~4.4]{WeibelEtaleChern} gives
\[
 i\bigl(c_{i,1}(N_{L/\Qp}y)
  -\operatorname{cor}_{L/\Qp}c_{i,1}(y)\bigr)=0
 \qquad (y\in K_d(L;\Z/p)).
\]
Because $i$ is a unit modulo $p$, the two maps agree.  On
$H^1(\Qp,\mu_p^{\otimes i})$, restriction followed by corestriction
is multiplication by $p-1$.  Given $h$ in this group, choose $y$ with
$c_{i,1}(y)=(p-1)^{-1}\operatorname{res}_{L/\Qp}h$.
Then $c_{i,1}(N_{L/\Qp}y)=h$, proving surjectivity over $\Qp$.
The compatible edges identify both the source and target of this
mod-$p$ map with $\Fp$.  It is the reduction of $c_i^{\mathrm S}$,
so that reduction is an isomorphism.  Both integral modules are
free of rank one; consequently $c_i^{\mathrm S}$ is an isomorphism.

It remains to identify this integral map on discrete classes.
All the coefficient identifications above are induced by completion
and localization, so the reduction of $\widehat x$ corresponds to the
ordinary reduction of the localized class $x$.
By \cite[Lemma~2.3]{WeibelEtaleChern}, the finite Chern map on this
reduction is evaluation of the same universal etale Chern class on
the integral Hurewicz image of $x$.  The proof of Proposition~2.8
there makes these evaluation maps compatible already in the
coefficient triangles.  For a fixed $x$, its Hurewicz image is
represented by a finite bar cycle at a finite matrix rank.  Evaluating
the compatible integral universal class on that cycle and reducing
therefore gives $c_{i,r}^{\mathrm S}(\widehat x\bmod p^r)$ at every
$r$.  The isomorphism $M_i=\varprojlim_r H^1(\Qp,\mu_{p^r}^{\otimes i})$
makes the integral evaluation uniquely equal to
$c_i^{\mathrm S}(\widehat x)$.

The regulator in \cite[Section~1.3]{HuberKings} evaluates the
universal standard-representation etale Chern class with $\Qp(i)$
coefficients on the same Hurewicz image.  It is the coefficient image
of the usual integral class just used: these Chern classes are
normalized on line bundles and characterized by the splitting
principle.  Hence, with the common suspension convention,
\[
 r_p(x)=c_i^{\mathrm S}(\widehat x)\otimes1
 \quad\text{in }H^1(\Qp,\Qp(i))
 \qquad (x\in K_d(\Zp)).
\]
Here the target is $M_i[1/p]$.  Indeed a continuous map from a compact
power of $G_{\Qp}$ to $\Qp(i)$ has bounded image, so it lies in
$p^{-N}\Zp(i)$ for some $N$.  Thus rational continuous cochains are
the localization of integral continuous cochains, and localization is
exact.  This comparison evaluates a fixed finite bar cycle and uses
the inverse limit in the target; it makes no exchange of an unbounded
group-cohomology limit with rationalization.
\end{proof}

We use the Bloch--Kato exponential of
\cite[Definition~3.10 and Proposition~1.17]{BlochKatoTamagawa};
its specialization to $\Qp(i)$ and the isomorphism for $i>1$
are recalled in \cite[Section~1.3]{HuberKings}.

\begin{lemma}[The integral exponential lattice]
\label{gen:exponential-lattice}
For $2\leq i\leq p-2$, use the standard period basis to identify
$D_{\mathrm{dR}}(\Qp(i))$ with $\Qp$.  Then
\[
 \exp_{\mathrm{BK}}^{-1}
     H^1(\Qp,\Zp(i))=p^i\Zp.
\]
\end{lemma}

\begin{proof}
For $i\geq2$, the rational group $H^1(\Qp,\Qp(i))$ is crystalline
and is identified with $\Qp$ by the Bloch--Kato exponential.  The
integral group $H^1(\Qp,\Zp(i))$ is free of rank one.  Indeed local
Galois cohomology is finitely generated, its rational rank is one,
and its torsion vanishes because
$H^0(\Qp,(\Qp/\Zp)(i))=0$: the mod-$p$ cyclotomic character to the
$i$th power is nontrivial in the stated range.

We shall construct one primitive integral extension using
Fontaine--Laffaille theory over the absolutely unramified field
$\Qp$.  We use covariant filtered Frobenius conventions with weights
$-i$ and $0$.  Taking the filtered dual puts these weights in
$[0,i]\subset[0,p-2]$, where the original contravariant
Fontaine--Laffaille functor applies.  Its exact lattice construction
is \cite[Section~7.14 and Proposition~7.15(i)]{FontaineLaffaille};
its full faithfulness on the reductions modulo $p$ is
\cite[Theorem~6.1]{FontaineLaffaille}.  The latter applies since
the interval avoids filtration degree $p-1$.  The rational comparison
in \cite[Theorem~8.4(ii)]{FontaineLaffaille} identifies the functor on
the filtered dual with the original covariant crystalline
representation.  Thus strongly divisible extensions in our
conventions give integral crystalline extensions of $\Zp$ by
$\Zp(i)$, compatibly with reduction and rationalization.

The rank-one endpoint has the standard period scale.  Let
$C=\widehat{\overline{\Qp}}$, $R=\mathcal O_{C^\flat}$, and
$\theta:W(R)\to\mathcal O_C$.  We also write $\theta$ for the
residue map $B_{\mathrm{dR}}^+\to C$.  Choose
$\xi=[x_0]+p$ with $x_0^\sharp=-p$, so $\ker\theta=(\xi)$.
At $q=p$, the period ring and its filtered pieces are
\[
 \mathscr S=W(R)[\xi^p/p],\qquad
 \mathscr S^j=(\xi^j,\xi^p/p)\quad(0<j<p)
\]
by \cite[Section~2.5 and Lemma~5.4]{FontaineLaffaille}; write
$\widehat{\mathscr S}^{\,j}$ for their separated $p$-adic completions
as in Section~7.14 there.  Choose a generator $e_{\mathrm{cyc}}$ of
$\Zp(1)$, represented by a primitive compatible root system
$\epsilon=(1,\zeta_p,\zeta_{p^2},\ldots)$, and put $\mu=[\epsilon]-1=\xi a$ with $a\in W(R)$.
The logarithm $t_{\mathrm{cyc}}=\log[\epsilon]$ belongs to
$\widehat{\mathscr S}^{\,1}$.  Indeed, if $k=pm+r$ with $0\leq r<p$,
\[
 \frac{\mu^k}{k}
 =\frac{p^m}{k}\,\xi^r a^k\left(\frac{\xi^p}{p}\right)^m.
\]
Here $m-v_p(k)\geq0$ and tends to infinity with $k$.  Each term lies
in $\mathscr S^1$ (when $r=0$, $m\geq1$), and the logarithm series
converges in its $p$-adic completion.  The displayed ideals give
$(\mathscr S^1)^i\subseteq\mathscr S^i$, while
$\phi(t_{\mathrm{cyc}})=p t_{\mathrm{cyc}}$
\cite[Section~4.11]{FontaineBarsottiTate}.
Here the completed divided Frobenius is
$\phi_i=p^{-i}\phi:\widehat{\mathscr S}^{\,i}\to\widehat{\mathscr S}$.
Thus
$t_{\mathrm{cyc}}^i\in\widehat{\mathscr S}^{\,i}$ and
$\phi_i(t_{\mathrm{cyc}}^i)=t_{\mathrm{cyc}}^i$.

The same ideal formula shows that
$x\mapsto\theta(x/\xi^i)$ sends $\mathscr S^i$ to $\mathcal O_C$:
the generator $\xi^p/p$ contributes zero because $i<p$.
This map extends $p$-adically to $\widehat{\mathscr S}^{\,i}$ and,
under the filtered period embedding in the proof of
\cite[Theorem~8.4(ii)]{FontaineLaffaille}, is the same residue map.
For the primitive period just chosen, the proof of
\cite[Proposition~2.17, pp.~543--544]{FontaineBarsottiTate} computes
$v_p\theta(t_{\mathrm{cyc}}/\xi)=1/(p-1)$.  Consequently
\[
 v_p\theta\left(\frac{t_{\mathrm{cyc}}^i/p}{\xi^i}\right)
 =\frac{i}{p-1}-1<0,
\]
so $t_{\mathrm{cyc}}^i/p\notin\widehat{\mathscr S}^{\,i}$.
For the positive rank-one object of weight $i$ with $\phi_i(f_i)=f_i$, its
realization is
$L_i=\{s\in\widehat{\mathscr S}^{\,i}:\phi_i(s)=s\}$.
Rank preservation and rational comparison make $L_i$ a rank-one
lattice in $\Qp t_{\mathrm{cyc}}^i$.  It contains
$t_{\mathrm{cyc}}^i$ but not $t_{\mathrm{cyc}}^i/p$, hence
$L_i=\Zp t_{\mathrm{cyc}}^i$.  The unit $1$ is primitive at weight
zero since $\theta(\widehat{\mathscr S})\subseteq\mathcal O_C$.
The corresponding covariant basis is therefore the
standard period basis
$e_i=t_{\mathrm{cyc}}^{-i}\otimes e_{\mathrm{cyc}}^{\otimes i}$.

Here is the lattice calculation in those conventions.  Let $e_i$ be
the standard basis of the rank-one filtered Frobenius object for
$\Qp(i)$, with $\phi(e_i)=p^{-i}e_i$, and choose an integral lift
$e_0$ of the quotient basis.  Write
\[
 \operatorname{Fil}^0=\Zp(e_0+x e_i),
 \qquad \phi(e_0)=e_0+y e_i.
\]
Strong divisibility is exactly
\begin{equation}\label{gen:fl-conditions}
 x\in\Zp,\qquad z:=y+p^{-i}x\in\Zp.
\end{equation}
Indeed the second condition says that $\phi(\operatorname{Fil}^0)$
is integral.  The condition in filtration $-i$ follows because
$p^iy=p^iz-x$ is integral, and generation holds because
$p^i\phi(e_i)=e_i$ while the degree-zero filtered generator maps to
a lift of $e_0$.  These are all the conditions for the two-step
filtered object.  Subtracting the rational Frobenius splitting gives
the exponential parameter
\begin{equation}\label{gen:fl-parameter}
 x-\frac{y}{1-p^{-i}}
   =\frac{x-z}{1-p^{-i}}.
\end{equation}
As $x,z$ range independently over $\Zp$, this ranges exactly over
$p^i\Zp$, since $p^i-1$ is a unit.  Replacing $e_0$ by
$e_0+b e_i$, $b\in\Zp$, replaces both $x$ and $z$ by their
differences with $b$, so \eqref{gen:fl-parameter} is independent of
this change of integral lift.

To see that the resulting lattice is the whole integral Galois
cohomology lattice, take $x=1$ and $z=0$.  Its extension modulo $p$
does not split in the filtered category.  This can be checked directly
in the standard positive interval: if $f_i,f_0$ are dual to $e_i,e_0$,
the dual lattice has
\[
 \phi(f_0)=f_0,\qquad \phi(f_i)=p^if_i+f_0,\qquad
 \operatorname{Fil}^{j}=\Zp(f_i-f_0)\quad(1\leq j\leq i).
\]
Its divided Frobenius satisfies $\phi_i(f_i-f_0)=f_i$.
A filtered section of the weight-$i$ quotient modulo $p$ must lift
its basis to $f_i-f_0$; commuting with $\phi_i$ would require its
image to be $f_i-f_0$ instead of $f_i$.  Thus no such section exists.
In the original coordinates this is the invariant $x-z=1$ under
all changes of lift, even after reduction.
Full faithfulness on finite objects implies that the corresponding
Galois extension modulo $p$ does not split either: a Galois splitting
would lift to a filtered splitting.  Its class in
$H^1(\Qp,\Zp(i))$ is consequently not divisible by $p$, and hence
generates that rank-one free module.  Its exponential parameter is
$p^i/(p^i-1)$ by \eqref{gen:fl-parameter}.  This proves
$\exp_{\mathrm{BK}}^{-1}H^1(\Qp,\Zp(i))=p^i\Zp$ using the
exact construction, finite full faithfulness, and rational
comparison just cited.
\end{proof}

The integral lattice fixes the scale of a generator of local Galois
cohomology.  To compare that scale with the continuous primitive
$b_i$, we will evaluate both classes on a discrete algebraic $K$-class.
The following lemma supplies a class on which this comparison detects
a nonzero scalar.

\begin{lemma}[A detected discrete regulator class]
\label{gen:detected-regulator}
For $2\leq i\leq p-2$, the Soul\'e regulator
\[
 r_p:K_{2i-1}(\Zp)\otimes_{\Z}\Qp
       \longrightarrow H^1(\Qp,\Qp(i))
\]
is nonzero.
\end{lemma}

\begin{proof}
Choose a nonidentity Teichmuller root of unity $\xi\in\Zp^\times$
and put $F=\Q(\xi)$.  The symbol $[\xi]_i$ lies in degree one of
de Jeu's rational weight-$i$ polylogarithmic complex over $F$
\cite{DeJeuWedge}, in the presentation of
\cite[p.~868]{BesserDeJeu}.  Its differential is
$[\xi]_{i-1}\otimes\xi$ for $i>2$ and $(1-\xi)\wedge\xi$ for $i=2$;
both vanish because $\xi$ is torsion in $F^\times$, hence zero in
$F^\times\otimes_\Z\Q$.  The first map in
\cite[Theorem~1.12]{BesserDeJeu} sends the resulting degree-one
cohomology class to $K_{2i-1}(F)\otimes_\Z\Q$.  For the chosen
embedding $F\to\Qp$, let $\mathcal O$ be the localization of the
integers of $F$ at the corresponding prime above $p$.  Both $\xi$ and
$1-\xi$ are units there.  The same theorem sends the rational
$K$-class through the localization isomorphism and the map to local
algebraic $K$-theory before applying the syntomic regulator.  This
number-field case requires no conjectural weight assumption; see
\cite[p.~871, the paragraph preceding Theorem~1.10, and Remark~1.13]
{BesserDeJeu}.
The regulator here is the syntomic Chern character in the normalized
coordinate of \cite[Definition~4.6 and Appendix~A]{BesserDeJeu}.
Its value is $\pm(i-1)!\operatorname{Li}_i(\xi)$, since the logarithmic
correction terms vanish at $\xi$.

There is some such $\xi$ for which this value is nonzero.  To see
this directly, use the Coleman measure identity
\begin{equation}\label{gen:coleman-measure}
 (1-p^{-i})\operatorname{Li}_i(\xi)
  =\lim_{r\to\infty}
    \frac{\displaystyle
       \sum_{\substack{1\leq b<p^r\\p\nmid b}}\xi^b b^{-i}}
         {1-\xi^{p^r}},
 \qquad \xi^p=\xi\neq1;
\end{equation}
see \cite[Proposition~4.8 and Corollary~4.9]{BesserDeJeuAlgorithm},
where the proof of Proposition~4.8 attributes the measure formula to
\cite[Lemma~7.2]{Coleman}.
This is integration of $x^{-i}$ on
$\Zp^\times$ against the bounded measure whose mass on
$b+p^r\Zp$ is $\xi^b/(1-\xi^{p^r})$.
The masses are compatible under refinement by the geometric-series
identity, so their integral modulo $p$ is computed at level one.
This finite-polylogarithm reduction is the one developed in
\cite[Proposition~2.1 and Corollary~2.2]{BesserFinitePolylog}.
Its numerator is the finite polylogarithm
\[
 P_i(X)=\sum_{b=1}^{p-1} b^{-i}X^b\in\Fp[X].
\]
This polynomial is nonzero, has degree $p-1$, and vanishes at both
$0$ and $1$; the latter uses $1\leq i\leq p-2$.
It cannot vanish at every element of $\Fp$, since that would give
at least $p$ distinct roots.  Some
$\overline\xi\in\Fp^\times\setminus\{1\}$ therefore has
$P_i(\overline\xi)\neq0$, and its Teichmuller lift makes the
right side of \eqref{gen:coleman-measure} nonzero.  The denominator
is a unit.  Since $1-p^{-i}\ne0$, this proves
$\operatorname{Li}_i(\xi)\ne0$, and hence the normalized
Chern-character regulator above is nonzero.

On Hurewicz images from positive algebraic $K$-groups, evaluation of
the usual Chern-character component of degree $i$ is
$(-1)^{i-1}/(i-1)!$ times evaluation of $c_i$; see
\cite[the remark after Definition~4.23 and Section~5.9]{TammeRelativeChern}.
The comparison from the syntomic theory used by Huber--Kings to the
modified theory above preserves Chern classes, by
\cite[Section~4]{BesserDeJeu} and
\cite[Propositions~2.2.7 and~2.2.9]{HuberKings}.  Because the normalized coordinate
is an isomorphism, its nonzero character value implies a nonzero
syntomic Chern value in the Huber--Kings theory.  Their isomorphism
$\eta:H^1_{\mathrm{syn}}(\Zp,i)\xrightarrow{\sim}\Qp$ is defined in
\cite[Definition~2.2.5]{HuberKings}.  By Proposition~2.3.4 there,
evaluation of this Chern class followed by $\eta$ and the Bloch--Kato
exponential is the etale regulator $r_p$.  The exponential is an
isomorphism for $i>1$ by Section~1.3 there.  Thus the
preceding construction supplies a discrete rational algebraic
$K$-theory class $x$ with $r_p(x)\ne0$.  If it is initially regarded
over $\Qp$, it has a preimage in $K_{2i-1}(\Zp)\otimes\Qp$ by
\cite[Remark~1.3.1]{HuberKings}, since $i>1$; the regulator extends $\Qp$-linearly.
\end{proof}

We can now determine the normalization of the classes selected by
the fixed local $K$-theory coordinates.  The argument needs the
integral exponential lattice and a single detected evaluation; the
two preceding lemmas supply these separately.

\begin{lemma}[Normalization of the selected primitives]
\label{lem:regulator-lattice}
With the choices of Lemma~\ref{gen:local-k-coordinates}, there are
units $\varepsilon_i\in\Zp^\times$, independent of the matrix rank,
such that
\begin{equation}\label{gen:primitive-normalization}
 \alpha_{a,i}\otimes\Qp=\varepsilon_i p^{-i}b_i
 \qquad (2\leq i\leq a\leq p-2).
\end{equation}
The same formula holds for $i=1$, with a unit
$\varepsilon_1$ determined by the divided logarithm.
\end{lemma}

\begin{proof}
For $i\geq2$, put
$g_i=(f_i)_*$ and define
\[
 t_i=\exp_{\mathrm{BK}}^{-1}
       \bigl(c_i^{\mathrm S}(g_i^{-1}(1))\bigr).
\]
Lemma~\ref{gen:exponential-lattice} and
Lemma~\ref{gen:soule-coordinate} give $t_i\in p^i\Zp^\times$.
Write $\alpha_i$ for the stable compatible family of leading classes,
and let $\alpha_i^\delta$ and $b_i^\delta$ denote the stable discrete
restrictions of the compatible continuous classes.  The bottom
Postnikov composite
\[
 K(\Zp)^\wedge_p\xrightarrow{f_i}\Sigma^{2i-1}\ell_p
       \longrightarrow\Sigma^{2i-1}H\Zp
\]
induces $g_i$ on $\pi_{2i-1}$.  Its fundamental ordinary class is
the leading class $\alpha_i$.  Lemma~\ref{gen:k-continuity} and
naturality of Hurewicz evaluation give
\[
 \langle\alpha_i^\delta,\operatorname{hur}(x)\rangle
       =g_i(\widehat x)
 \qquad (x\in K_{2i-1}(\Zp)),
\]
and the same identity extends $\Qp$-linearly to
$K_{2i-1}(\Zp)\otimes_{\Z}\Qp$.  The fundamental
class of the Eilenberg--Mac Lane target evaluates as the identity on
its homotopy group, so no additional normalization factor occurs.
By \cite[Definition~1.2.3 and Theorem~1.3.2]{HuberKings} and
Lemma~\ref{gen:soule-coordinate},
\[
 \langle b_i^\delta,\operatorname{hur}(x)\rangle
   =\exp_{\mathrm{BK}}^{-1}r_p(x)
   =t_i g_i(\widehat x)
   =t_i\langle\alpha_i^\delta,\operatorname{hur}(x)\rangle.
\]
Choose $x$ as in Lemma~\ref{gen:detected-regulator}.  By
Lemma~\ref{gen:rational-components}, write the stable continuous
primitive as $\alpha_i\otimes\Qp=\lambda_i b_i$.  The evaluation
identities above are nonzero on $x$ and give
$\lambda_i=t_i^{-1}=\varepsilon_i p^{-i}$ for a unit
$\varepsilon_i\in\Zp^\times$.  Both $g_i$ and $c_i^{\mathrm S}$
were fixed independently of matrix rank, so this is the same unit
at every rank.  This uses one detected discrete evaluation and
requires no density assertion about discrete local algebraic
$K$-theory.

Finally, $b_1$ corresponds under the Lazard comparison to the ordinary
determinant logarithm.  Since $\log(1+p)$ has valuation one, the
normalization in Lemma~\ref{gen:sphere-representatives} gives
$\alpha_{a,1}=\varepsilon_1p^{-1}b_1$ with
$\varepsilon_1\in\Zp^\times$.
\end{proof}

\subsection{Top volume and the integral product}
\label{gen:volume-section}

Let $d=a^2$.  Order the $a^2$ matrix entries once and for all, and
write
\[
 \omega_a=\bigwedge_{r,s}(g^{-1}dg)_{rs}
          =\det(g)^{-a}\bigwedge_{r,s}dg_{rs}
\]
for the corresponding invariant algebraic top differential.  Its
value at the identity is the raw entry volume on
$\mathfrak{gl}_a(\Qp)$.  The normalization just proved gives the
product $\alpha_{a,1}\cdots\alpha_{a,a}$ a factor
$p^{-a(a+1)/2}$ relative to $b_1\cdots b_a$.  We now show that
this is exactly the integral top-cohomology lattice of $G_a$.
This comparison will prove that the selected product generates the
integral top cohomology.

\begin{lemma}[The top integral lattice]\label{lem:top-volume}
Assume $a\leq p-2$.
\begin{enumerate}
\item In rational Lie algebra cohomology,
\[
 b_1\wedge\cdots\wedge b_a=u_a\omega_a,
 \qquad u_a\in\Z_{(p)}^\times.
\]
\item For all sufficiently large integers $r$, the integral top
cohomology lattice of
$U_r=1+p^rM_a(\Zp)$ corresponds under the rational Lazard
comparison to
\[
 \Zp\,p^{-rd}\omega_a.
\]
\item The top integral lattice of $G_a$ corresponds to
\[
 \Zp\,p^{-a(a+1)/2}\omega_a.
\]
\end{enumerate}
\end{lemma}

\begin{proof}
For the first assertion, we use the Chern-class comparison and
volume calculation of Huber--Soergel
\cite[Propositions~1.3, 3.3, and 4.1]{HuberSoergel}.
The compatibility with suspension in their Corollary~3.4 identifies
precisely the primitives $b_i$ fixed above, rather than arbitrary
rational generators.  On the compact real form $U(a)$ of
$\GL_a(\mathbb C)$, these transgressed Chern classes, divided by their
Tate factors $(2\pi\sqrt{-1})^i$, are integral odd generators, and
their product integrates to $1$ up to orientation.
Successive last-column projections
$U(j)\to S^{2j-1}$ compute the entry-volume integral as
\[
 \int_{U(a)}\omega_a
   =\pm\frac{(2\pi\sqrt{-1})^{a(a+1)/2}}
                 {\prod_{j=1}^a(j-1)!}.
\]
One obtains this formula by using the sphere volume
$2\pi^j/(j-1)!$ at the $j$th step; each complex off-diagonal tangent
entry contributes its $2\sqrt{-1}$ factor and the imaginary diagonal
entry its $\sqrt{-1}$ factor.  Comparing with the product of the
transgressions gives exactly $u_a=\pm\prod_{j=1}^a(j-1)!$ in the
chosen normalization.  Since $a\leq p-2$, this coefficient is a
$p$-unit.  In particular this calculation does not introduce
the possibly nonunit factor $(a^2)!$ into the primitive product.

We give a direct integral verification of the second assertion that
also handles that latter factorial.  Put
$e=v_p(d!)$, and take $r$ sufficiently large; for example $r\geq d+1$
is sufficient for the estimates below.  In the affine coordinates
$g=1+p^rY$ on $U_r$, the normalized form is
\[
 p^{-rd}\omega_a
       =\det(1+p^rY)^{-a}\bigwedge_{r',s'}dY_{r's'}.
\]
Integrate this form formally, term by term, over affine straight
simplices with vertices in $U_r$.  The integral of a monomial of
total degree $h$ on a $d$-simplex has denominator dividing
$(h+d)!$.  The expansion
\[
 \det(1+p^rY)^{-a}
      =1+\sum_{h\geq1}p^{rh}Q_h(Y)
\]
has integral homogeneous polynomials $Q_h$.  Hence the resulting
series of simplex integrals converges $p$-adically.  Left
multiplication on matrices is affine and preserves $\omega_a$.
Formal Stokes and the faces of a simplex therefore give a homogeneous
group cocycle.  In inhomogeneous notation call it $C_r$, evaluated
on the vertices
\[
 1,\ g_1,\ g_1g_2,\ \ldots,\ g_1\cdots g_d.
\]
It is a continuous rational analytic cocycle representing the
normalized form under differentiation.  Its leading term is
\[
 \frac{1}{d!}\det(X_1,\ldots,X_d),
 \qquad X_j=(g_j-1)/p^r.
\]
Differentiation alternates in the $d$ variables and cancels precisely
this $d!$, so the corresponding Lie cohomology class is
$p^{-rd}\omega_a$ with the stated normalization.  This description
of the rational Lazard map agrees with the analytic differentiation
comparison in \cite[Section~4]{HuberKings}; an all-degree integral
comparison for uniform groups is also supplied by
\cite[Theorem~3.3.3, Proposition~4.2.4, and Remark~4.2.5]
{HuberKingsNaumann}.

The following reduction proves the integral assertion directly.
Multiply $C_r$ by $d!$ and work modulo $p^{e+1}$.  Every term with
$h\geq1$ has valuation at least
\[
 e+rh-v_p((h+d)!)\geq e+1
\]
for the chosen $r$.  For instance
$v_p((h+d)!)\leq(h+d)/(p-1)$ proves this inequality when
$r\geq d+1$.  Furthermore successive-product vertices differ from
their additive partial sums by multiples of $p^r$.  Since
$r\geq e+1$, these differences do not affect the reduction.
Writing $A_e=\Z/p^{e+1}$, the entry functions
\[
 \chi_j:U_r\longrightarrow A_e,
 \qquad \chi_j(g)=((g-1)/p^r)_j\pmod {p^{e+1}},
 \quad 1\leq j\leq d,
\]
are genuine group homomorphisms.  Indeed the additional product term
is divisible by $p^r$.  We consequently obtain
\begin{equation}\label{gen:factorial-reduction}
 [d!C_r]
  =d![\chi_1\smile\cdots\smile\chi_d]
   \quad\hbox{in }H^d_{\mathrm{cts}}(U_r,A_e).
\end{equation}
Here the determinant cocycle is the alternating sum of the ordered
cups.  Graded commutativity makes each term equal to the ordered cup
after its permutation sign is included, giving the factor $d!$.

The group $U_r$ is uniform, and the reductions of the $\chi_j$ are its
standard degree-one basis.  Their ordered cup is the nonzero top
generator in mod-$p$ cohomology.  Orientable compact-group duality
gives
\[
 H^d_{\mathrm{cts}}(U_r,\Zp)=\Zp,
 \qquad H^{d+1}_{\mathrm{cts}}(U_r,\Zp)=0.
\]
Reduction thus identifies $H^d(U_r,A_e)$ with
$H^d(U_r,\Zp)/p^{e+1}$.  The ordered cup in
\eqref{gen:factorial-reduction} is a unit in this cyclic module.
The integral cocycle $d!C_r$ therefore represents a class of
\emph{exactly} valuation $e$ in $H^d(U_r,\Zp)$.  Dividing its
rational class by $d!$ proves that $[C_r]$ is an integral unit
generator.  This argument does not divide by $d!$ inside $A_e$;
the valuation is first determined by reduction and only then divided
in rational cohomology.  It proves the second assertion even when
$p\mid d!$.

For the last assertion, the determinant of the adjoint action of
$G_a$ is one, so its orientation is trivial.  Top restriction from
$G_a$ to $U_r$ multiplies an integral orientation generator by
$[G_a:U_r]$, up to a unit; this follows equally by dualizing degree-zero
corestriction or by the restriction--corestriction identity and
oriented duality.  Its index has valuation
\begin{align*}
 v_p[G_a:U_r]
  &=(r-1)a^2+v_p|\GL_a(\Fp)|\\
  &=(r-1)a^2+\frac{a(a-1)}2
   =ra^2-\frac{a(a+1)}2.
\end{align*}
The rational Lie algebra is unchanged under restriction.  Combining
this index with the lattice $p^{-ra^2}\omega_a$ gives
$p^{-a(a+1)/2}\omega_a$ for the full group.
\end{proof}

\subsection{Products, block restrictions, and transport}
\label{gen:products-section}

\begin{proof}[Proof of Theorem~\ref{thm:stable-generators}]
By Lemmas~\ref{lem:regulator-lattice} and~\ref{lem:top-volume}, the
product of the leading integral classes has rational image
\[
 \alpha_{a,1}\cdots\alpha_{a,a}
   =\left(\prod_{i=1}^a\varepsilon_i\right)
      p^{-a(a+1)/2}\,b_1\cdots b_a
\]
and is an integral unit generator in top degree $a^2$.  In particular
its mod-$p$ reduction is nonzero.

All the classes have odd degree, and $2$ is invertible in $\Zp$, so
they define a map from the indicated exterior algebra.  If the graded kernel of the mod-$p$ map contained a nonzero
homogeneous element, the perfect top-degree pairing in the source
exterior algebra would multiply it to a nonzero top class in the
kernel.  The nonvanishing top image therefore makes this map
injective.
Theorem~\ref{thm:constant-cohomology} gives the same graded dimensions
on both sides, so it is an isomorphism.  Integral torsion-freeness,
finite generation in each degree, and reduction modulo $p$ then show
that the integral map is also an isomorphism.  This proves the
generator assertions, with actual leading classes rather than a
choice of abstract exterior generators.

The definitions of $\kappa_a$ and $f_i$ make their leading classes
compatible with every upper-left block restriction.  For $i>b$, the
rational primitive restricts to zero at rank $b$.  Integral
torsion-freeness at that rank makes the integral restriction zero as
well.  For $i\leq b$ it is exactly $\alpha_{b,i}$, since the
projection $f_i$ and its free-coordinate normalization were fixed
independently of rank.  To choose compatible sphere representatives,
first use Lemma~\ref{gen:sphere-representatives} at rank $N$ and then
restrict those chosen representatives to every smaller rank.  Their
images in $\ell_p$ and their leading integral classes are the ones
already constructed there.  No simultaneous canonical choice of
nullhomotopies is needed.

Finally apply Proposition~\ref{prop:modification-transport} to these
rank-$N$ sphere classes.  That proposition transports continuous
stable lifts through the common modification stack and preserves
their mod-$p$ Hurewicz images.  It also identifies their action after
choosing the connected quotient frame with the actual upper-left
matrix-block restriction.  Thus the transported classes $u_i$ have
the leading classes $\xi_i$ stated in the theorem, and for $i\leq m$
their action on the extension space is through
$\overline\alpha_{m,i}$.  This proves all the asserted
compatibilities.
\end{proof}
